% Part: lambda-calculus
% Chapter: introduction
% Section: primitive-recursion

\documentclass[../../../include/open-logic-section]{subfiles}

\begin{document}

\olfileid[fr]{lam}{int}{pr}
\olsection{Les fonctions \texorpdfstring{$\lambd$}{λ}-définissables et la récurrence primitive}

Pour la récurrence primitive, un travail supplémentaire
s'impose. Nous procéderons par étapes. Supposons donnés
des termes $G'$ et~$H'$ qui définissent respectivement
les fonctions $g$ et~$h$. Nous cherchons un terme~$F$
qui définisse la fonction~$f$ donnée par
\begin{align*}
f(0, \vec z) & = g(\vec z) \\
f(x+1, \vec z) & = h(x, f(x,\vec z), \vec z).
\end{align*}
Il suffit donc de construire~$F$ tel que, pour tout
entier naturel~$n$,
\begin{align*}
F(\num{0}, \vec z) & \equiv G'(\vec z) \\
F(\num{n+1}, \vec z) & \equiv
  H'(\num{n}, F(\num{n}, \vec z), \vec z).
\end{align*}
Ici et dans la suite, $\equiv$ désigne
l'équivalence~$\beta$. Lorsque nous remplaçons
$\vec z$ par des numéraux $\num{\vec m}$, le fait
que $G'$ et~$H'$ définissent $g$ et~$h$ donne alors
la valeur attendue, pourvu que les calculs
intermédiaires soient définis.

Nous pouvons absorber les paramètres supplémentaires
$\vec z$ dans des abstractions. Posons
\begin{align*}
G & = \lambd[\vec z][G'(\vec z)], \\
H(u,v) & = \lambd[\vec z][H'(u,v(\vec z),\vec z)].
\end{align*}
Il suffit alors de trouver~$F$ satisfaisant
\begin{align*}
F(\num{0}) & \equiv G, \\
F(\num{n+1}) & \equiv H(\num{n},F(\num{n}))
\end{align*}
pour tout~$n$. Cette curryfication nous dispense
de manipuler explicitement les paramètres~$\vec z$
à chaque étape.

Avant de définir~$F$, construisons un terme qui
représente les couples ordonnés.

\begin{lem}
Il existe un lambda-terme~$D$ tel que, pour tous
termes $M$ et~$N$,
$D(M,N)(\num{0})\red M$ et
$D(M,N)(\num{1})\red N$.
\end{lem}

\begin{proof}
Définissons d'abord le terme~$K$ par
\[
K(y)=\lambd[x][y].
\]
Autrement dit, $K=\lambd[y][\lambd[x][y]]$.
Pour tout~$M$, $K(M)$ est donc la fonction
constante qui renvoie~$M$ quel que soit son
argument.

Définissons maintenant $D(x,y,z)$ par
$D(x,y,z)=z(K(y))x$. Alors
\begin{align*}
D(M,N,\num{0}) & \red \num{0}(K(N))M \red M,
  \quad\text{et}\\
D(M,N,\num{1}) & \red \num{1}(K(N))M
  \red K(N)M \red N,
\end{align*}
comme voulu.
\end{proof}

Le terme $D(M,N)$ représente donc le couple
$\tuple{M,N}$. Si $P$ représente un tel couple,
$P(\num{0})$ et $P(\num{1})$ en donnent les
projections gauche et droite. Nous noterons
celles-ci $(P)_0$ et $(P)_1$.

\begin{lem}
Les fonctions totales $\lambd$-définissables sont
fermées par récurrence primitive.
\end{lem}

\begin{proof}
Étant donnés des termes $G$ et~$H$, nous cherchons
un terme~$F$ tel que
\begin{align*}
F(\num{0}) & \equiv G, \\
F(\num{n+1}) & \equiv H(\num{n},F(\num{n}))
\end{align*}
pour chaque entier naturel~$n$. L'idée est de
calculer, en utilisant les numéraux comme
itérateurs, la suite de \emph{couples}
\[
\tuple{\num{0},F(\num{0})},
\tuple{\num{1},F(\num{1})},\dots.
\]
Le premier couple est $\tuple{\num{0},G}$.
À partir de $\tuple{\num{n},F(\num{n})}$,
le couple suivant doit être équivalent à
$\tuple{\num{n+1},H(\num{n},F(\num{n}))}$.
Nous allons définir un terme~$T$ qui effectue
cette transition.

Posons
\[
T(u)=\tuple{\fn{Succ}((u)_0),H((u)_0,(u)_1)}.
\]
Pour tout~$n$ et tout terme~$M$, on vérifie que
\[
T(\tuple{\num{n},M})
\red\tuple{\num{n+1},H(\num{n},M)}.
\]
Étant donnés $G$ et~$H$, définissons~$F(u)$ par
\[
F(u)=(u(T,\tuple{\num{0},G}))_1.
\]
Sur l'entrée~$\num{n}$, $F$ applique donc~$T$
$n$~fois au couple $\tuple{\num{0},G}$, puis
renvoie la seconde composante. Au départ,
\begin{enumerate}
\item $\num{0}(T,\tuple{\num{0},G})
  \equiv\tuple{\num{0},G}$;
\item $F(\num{0})\equiv G$.
\end{enumerate}
Par récurrence sur~$n$, montrons les deux
relations suivantes :
\begin{enumerate}
\item $\num{n+1}(T,\tuple{\num{0},G})
  \equiv\tuple{\num{n+1},F(\num{n+1})}$;
\item $F(\num{n+1})
  \equiv H(\num{n},F(\num{n}))$.
\end{enumerate}
Pour la seconde, nous avons
\begin{align*}
F(\num{n+1})
&\red(\num{n+1}(T,\tuple{\num{0},G}))_1 \\
&\equiv(T(\num{n}(T,\tuple{\num{0},G})))_1 \\
&\equiv(T(\tuple{\num{n},F(\num{n})}))_1 \\
&\equiv(\tuple{\num{n+1},H(\num{n},F(\num{n}))})_1 \\
&\equiv H(\num{n},F(\num{n})).
\end{align*}
Le remplacement de l'itération par le couple utilise
l'hypothèse de récurrence. Pour la première
relation, nous avons
\begin{align*}
\num{n+1}(T,\tuple{\num{0},G})
&\equiv T(\num{n}(T,\tuple{\num{0},G})) \\
&\equiv T(\tuple{\num{n},F(\num{n})}) \\
&\equiv\tuple{\num{n+1},H(\num{n},F(\num{n}))} \\
&\equiv\tuple{\num{n+1},F(\num{n+1})}.
\end{align*}
La dernière ligne utilise la seconde relation.
Nous avons ainsi établi $F(\num{0})\equiv G$
et, pour tout~$n$,
$F(\num{n+1})\equiv H(\num{n},F(\num{n}))$,
comme souhaité.
\end{proof}

La construction ci-dessus donne la valeur attendue dès que tous
les calculs intermédiaires sont définis. Avec des fonctions
partielles, la récurrence stricte peut être indéfinie même si le
terme obtenu possède une forme normale. Si $g$ est partout
indéfinie et $h$ constante zéro, $f(1,\vec z)$ est indéfinie.
Pourtant, le terme~$H$ de la construction peut effacer le résultat
de $F(\num{0},\vec z)$ et renvoyer~$\num{0}$.
Pour la fermeture de la classe des fonctions partielles
$\lambd$-définissables, on utilise leur caractérisation par les
fonctions récursives partielles et la clôture de cette dernière
classe par récurrence primitive.

\end{document}
